\documentclass[12pt,a4paper]{article}
\usepackage[utf8]{inputenc}
\usepackage[english]{babel}
\usepackage{graphicx}
\usepackage{float}
\usepackage{amsmath}
\usepackage{amssymb}
\usepackage{hyperref}
\usepackage{listings}
\usepackage{xcolor}
\usepackage{geometry}
\usepackage{cite}
\usepackage{fancyhdr}
\usepackage{caption}
\usepackage{subcaption}
\usepackage{multirow}
\usepackage{array}

\geometry{a4paper, margin=1in}

% Code listing style
\lstdefinestyle{verilog}{
    language=Verilog,
    basicstyle=\ttfamily\small,
    keywordstyle=\color{blue}\bfseries,
    commentstyle=\color{green!60!black}\itshape,
    stringstyle=\color{red},
    numbers=left,
    numberstyle=\tiny\color{gray},
    stepnumber=1,
    numbersep=5pt,
    backgroundcolor=\color{gray!5},
    frame=single,
    breaklines=true,
    breakatwhitespace=true,
    tabsize=2,
    captionpos=b
}

% Header and footer
\pagestyle{fancy}
\fancyhf{}
\rhead{SAFAS - Hardware Scheduler}
\lhead{Mid-Term Project Report}
\cfoot{\thepage}

\title{
    \textbf{SAFAS: Secure And FAst Hardware Scheduler} \\
    \large{FPGA-Based Hardware Scheduler for Accelerating Task Scheduling in Multi-Core Systems} \\
    \vspace{0.5cm}
    \includegraphics[width=0.4\textwidth]{Logo-SAFAS.png}
}

\author{Mid-Term Project Report}
\date{\today}

\begin{document}

\maketitle
\thispagestyle{empty}

\newpage
\tableofcontents
\newpage

\begin{abstract}
This report presents a comprehensive analysis and implementation of SAFAS (Secure And FAst hardware Scheduler), an FPGA-based hardware scheduler designed for real-time task scheduling in multi-core embedded systems. With the proliferation of IoT devices and safety-critical applications, the need for secure and efficient task scheduling has become paramount. SAFAS addresses this challenge by implementing a hardware-based scheduler that operates independently from the operating system, providing enhanced security and performance through direct hardware control of processing units.

The system implements the Earliest Deadline First (EDF) scheduling algorithm in hardware, utilizing a sophisticated insertion-sort architecture that maintains tasks in deadline order. The design supports up to 16 processing cores, manages multiple priority queues (normal and security-critical tasks), and provides preemptive scheduling capabilities to ensure hard real-time deadlines are met. The architecture includes comprehensive fault detection, statistical monitoring, and security-aware task replication for critical operations.

This report details the architectural design, module-level implementation, verification methodology, and performance analysis of the SAFAS hardware scheduler, demonstrating its effectiveness for modern embedded systems requiring both high performance and security guarantees.
\end{abstract}

\newpage

\section{Introduction}

\subsection{Motivation}
The increasing complexity of modern embedded systems, particularly in the Internet of Things (IoT) domain, has created significant challenges for traditional software-based task scheduling approaches. Operating system schedulers often introduce unpredictable latencies, consume valuable processor cycles for scheduling decisions, and can be vulnerable to schedule-based attacks that exploit timing information to compromise system security.

Hardware schedulers offer a compelling alternative by:
\begin{itemize}
    \item \textbf{Reducing overhead:} Scheduling decisions occur in dedicated hardware, freeing CPU resources for application tasks
    \item \textbf{Guaranteeing determinism:} Hardware execution provides predictable, cycle-accurate timing
    \item \textbf{Enhancing security:} Separation of scheduling logic from software reduces attack surface
    \item \textbf{Enabling parallelism:} Direct hardware control of multiple cores allows efficient parallel task management
\end{itemize}

\subsection{Project Objectives}
The primary objectives of the SAFAS project are:
\begin{enumerate}
    \item Design and implement a hardware scheduler in Verilog HDL that can efficiently manage real-time tasks across multiple processor cores
    \item Implement the EDF (Earliest Deadline First) scheduling algorithm in hardware to optimize deadline adherence
    \item Provide security-aware scheduling with support for replicated execution of critical tasks
    \item Minimize hardware resource utilization while maintaining high performance
    \item Verify correctness through comprehensive simulation and testbench development
    \item Demonstrate feasibility for FPGA implementation targeting Xilinx Virtex family devices
\end{enumerate}

\subsection{System Overview}
SAFAS implements a complete hardware scheduling system consisting of:
\begin{itemize}
    \item \textbf{Main Queue:} Large buffer (256-512 tasks) for receiving tasks from the real-time operating system
    \item \textbf{Dual Ready Queues:} Separate queues for normal (64 tasks) and security-critical (8 tasks) operations
    \item \textbf{EDF Sorter:} Hardware insertion-sort network maintaining tasks in earliest-deadline-first order
    \item \textbf{Multi-Level Arbiter:} Fair processor selection for task assignment
    \item \textbf{Preemptive Dispatcher:} Two-situation scheduling logic supporting task preemption
    \item \textbf{Controller:} State machine managing timing and repair operations
    \item \textbf{Core Interface:} Connection to 16 processing cores via Network-on-Chip
\end{itemize}

\section{System Architecture}

\subsection{Overall Architecture}

Figure~\ref{fig:architecture} illustrates the complete SAFAS architecture, showing the interconnection of all major components.

\begin{figure}[H]
    \centering
    \includegraphics[width=0.95\textwidth]{MainArchitecture.jpg}
    \caption{SAFAS Hardware Scheduler Architecture showing task flow, control signals, and processor interface}
    \label{fig:architecture}
\end{figure}

The architecture follows a pipelined dataflow model:

\begin{enumerate}
    \item \textbf{Task Reception:} The real-time operating system validates and transmits tasks to the Main Queue
    \item \textbf{Queue Management:} Tasks are transferred from the Main Queue to the Scheduler's ready queues during active scheduling periods
    \item \textbf{EDF Sorting:} Both Ready Queue and Critical Queue maintain tasks in ascending deadline order using hardware insertion cells
    \item \textbf{Task Dispatching:} The scheduler dispatches tasks to processor cores based on two situations:
    \begin{itemize}
        \item \textbf{Situation 1:} Processor completes current task and requests new assignment
        \item \textbf{Situation 2:} Preemption when a new task has earlier deadline than currently running tasks
    \end{itemize}
    \item \textbf{Execution Monitoring:} Running tasks are continuously monitored for completion and deadline violations
    \item \textbf{Periodic Repair:} A repair mechanism identifies and removes tasks that have missed their deadlines
\end{enumerate}

\subsection{Task Format}

Tasks in SAFAS are represented as 41-42 bit data structures (depending on module):

\begin{table}[H]
\centering
\begin{tabular}{|c|c|l|}
\hline
\textbf{Bits} & \textbf{Width} & \textbf{Description} \\
\hline
[41] or [40] & 1 & VALID bit (1 = valid task, 0 = empty/invalid) \\
[40] (optional) & 1 & Type/Security bit (1 = critical, 0 = normal) \\
[39:32] & 8 & Task ID (unique identifier 0-255) \\
[31:16] & 16 & Relative Deadline (time units until deadline) \\
[15:0] & 16 & Execution Time (required processing cycles) \\
\hline
\end{tabular}
\caption{Task Data Structure Format}
\label{tab:task_format}
\end{table}

This compact representation allows efficient hardware processing while providing sufficient range for practical real-time systems (deadlines up to 65,535 time units).

\subsection{Timing and Control}

The Controller module (\texttt{4\_Controller.v}) implements a five-state state machine that orchestrates system timing:

\begin{table}[H]
\centering
\begin{tabular}{|l|c|p{6cm}|}
\hline
\textbf{State} & \textbf{Duration} & \textbf{Operations} \\
\hline
SCHEDULING & 64 cycles & Main scheduling operations: Situation 1 and Situation 2 checks, task dispatching \\
\hline
MARGIN-TIME & 64 cycles & Data propagation time for queue operations \\
\hline
SUBTRACT & 1 cycle & Simultaneous decrement of all task deadlines and execution times \\
\hline
IDLE & 65,407 cycles & Processor execution period (production) or 256 cycles (simulation) \\
\hline
MARGIN-TIME & 64 cycles & Preparation for next scheduling cycle \\
\hline
\end{tabular}
\caption{Controller State Machine Timing}
\label{tab:timing}
\end{table}

The total cycle period is approximately 65,536 cycles (16-bit counter), providing a predictable and deterministic scheduling rhythm. During production operation with typical FPGA clock frequencies of 200-250 MHz, this translates to scheduling periods of approximately 262-328 microseconds.

\section{Module-Level Implementation}

\subsection{Module 0\_REG.v - Parameterizable Register}

The basic building block for all synchronous storage in SAFAS is a parameterizable register module.

\begin{lstlisting}[style=verilog, caption={Register Module Interface}]
module REG #(parameter NUM = 16)
(
    input clk,
    input rst,
    input [NUM-1:0] IN,
    output reg [NUM-1:0] OUT
);

always @(posedge clk) begin
    if (rst)
        OUT <= {NUM{1'b0}};
    else
        OUT <= IN;
end

endmodule
\end{lstlisting}

\textbf{Key Features:}
\begin{itemize}
    \item Synchronous reset for predictable initialization
    \item Parameterizable width for flexible data sizes
    \item Single-cycle latency for pipeline stages
    \item Used throughout design for read-after-write hazard prevention
\end{itemize}

\subsection{Module 1\_Insertion\_cell.v - EDF Sorting Cell}

The insertion cell is the fundamental component implementing hardware-based EDF sorting. Each cell can store one task and compare incoming tasks against its stored value.

\textbf{Operating Principle:}

When a new task arrives with deadline $D_{new}$ and the cell contains a task with deadline $D_{stored}$:

\begin{equation}
    \text{keep\_stored} = (D_{stored} \leq D_{new}) \land \text{valid\_stored}
\end{equation}

\begin{equation}
    \text{output\_task} =
    \begin{cases}
        \text{stored\_task} & \text{if keep\_stored} \\
        \text{new\_task} & \text{otherwise}
    \end{cases}
\end{equation}

\begin{lstlisting}[style=verilog, caption={Insertion Cell Core Comparison Logic}]
wire comp = (data_reg[31:16] <= data_in[31:16]);

assign data_in_p = (comp & data_reg[W-1]) ? data_reg :
                   {1'b1, data_in[W-2:0]};
assign data_out_p = (comp & data_reg[W-1]) ?
                    {1'b1, data_in[W-2:0]} : data_reg;
\end{lstlisting}

\textbf{Subtraction Support:}

Each cell includes a \texttt{sub\_task} submodule that decrements both execution time and relative deadline during the SUBTRACT phase:

\begin{lstlisting}[style=verilog, caption={Deadline and Execution Decrement}]
module sub_task #(parameter W = 41)
(
    input [W-1:0] task_in,
    input subtract_en,
    input repair_period,
    output [W-1:0] task_out
);

wire check_miss = task_in[31:16] < task_in[15:0];

assign task_out[15:0] = subtract_en ?
                        (task_in[15:0] - 1) : task_in[15:0];
assign task_out[31:16] = subtract_en ?
                         (task_in[31:16] - 1) : task_in[31:16];
assign task_out[W-1] = (repair_period & check_miss) ?
                       1'b0 : task_in[W-1];

endmodule
\end{lstlisting}

The \texttt{check\_miss} signal detects deadline violations when remaining deadline becomes less than remaining execution time, making task completion impossible.

\subsection{Module 1\_Insertion\_block.v - Sorted Queue}

The insertion block chains $N$ insertion cells (typically $N=64$) to create a fully sorted queue maintaining EDF order.

\textbf{Architecture:}

\begin{figure}[H]
\centering
\begin{verbatim}
    data_in  → [Cell 0] → [Cell 1] → ... → [Cell N-1] → data_fail
                  ↓          ↓                  ↓
              data_out     (sorted queue)    (rejected)
\end{verbatim}
\caption{Insertion Block Cascaded Cell Architecture}
\end{figure}

\begin{lstlisting}[style=verilog, caption={Insertion Block Generate Statement}]
generate
    for (o = 0; o < N; o = o + 1) begin : IB
        if (o == 0) begin
            // First cell
            Insertion_cell #(.W(W)) Ins_cell(
                .clk(clk), .rst(rst),
                .data_in(data_in),
                .data_out(data_out),
                .data_in_p(data_IB[o]),
                // ... additional signals
            );
        end else if (o == N-1) begin
            // Last cell (overflow detection)
            Insertion_cell #(.W(W)) Ins_cell(
                .data_in_p(data_fail),
                .data_in(data_IB[o-1]),
                // ... additional signals
            );
        end else begin
            // Intermediate cells
            Insertion_cell #(.W(W)) Ins_cell(
                .data_in(data_IB[o-1]),
                .data_in_p(data_IB[o]),
                // ... additional signals
            );
        end
    end
endgenerate
\end{lstlisting}

\textbf{Performance Characteristics:}
\begin{itemize}
    \item \textbf{Insertion Time:} $O(N)$ cycles for pipeline fill, $O(1)$ per task after
    \item \textbf{Comparison Operations:} $N$ parallel comparisons per insertion
    \item \textbf{Critical Path:} Single cell comparison + multiplexer
    \item \textbf{Hardware Cost:} $N \times W$ flip-flops + $N$ comparators
\end{itemize}

\subsection{Module 2\_Arbiter.v - Hierarchical Processor Selection}

The arbiter implements fair processor selection when multiple cores simultaneously request new tasks (Situation 1).

\textbf{Architecture:} Four-level binary tree

\begin{figure}[H]
\centering
\begin{verbatim}
Level 1:  [P0,P1] [P2,P3] [P4,P5] [P6,P7] [P8,P9] [P10,P11] [P12,P13] [P14,P15]
            ↓       ↓       ↓       ↓       ↓        ↓         ↓         ↓
Level 2:    [A0,A1]       [A2,A3]         [A4,A5]           [A6,A7]
                ↓             ↓               ↓                 ↓
Level 3:         [B0,B1]                 [B2,B3]
                      ↓                       ↓
Level 4:                  [Selected Processor]
\end{verbatim}
\caption{Four-Level Hierarchical Arbiter Tree Structure}
\end{figure}

\textbf{Round-Robin Fairness:}

Each arbiter cell maintains a priority toggle that alternates between two inputs:

\begin{lstlisting}[style=verilog, caption={Fair Arbiter Cell Logic}]
always @(posedge clk) begin
    if (rst)
        priority_ <= 1'b0;
    else if (req[0] | req[1])
        priority_ <= !priority_;  // Toggle on each request
end

assign op1[0] = ~priority_ | ~req[1];
assign op1[1] = priority_ | ~req[0];
assign op2[0] = req[0] & op1[0];
assign op2[1] = req[1] & op1[1];
\end{lstlisting}

This ensures neither processor can monopolize the scheduler even under contention.

\textbf{Latency Analysis:}
\begin{itemize}
    \item \textbf{Combinational Path:} 4 levels of 2:1 multiplexers
    \item \textbf{Total Latency:} 4 clock cycles (one per level)
    \item \textbf{Grant Propagation:} Selected processor index propagates back through tree
\end{itemize}

\subsection{Module 3\_Scheduler\_partial.v - Core Scheduling Logic}

This module implements the heart of the SAFAS scheduling algorithm, managing dual queues and two-situation dispatching.

\subsubsection{Dual Queue System}

\textbf{Ready Queue (64 tasks):} Stores normal priority tasks
\begin{lstlisting}[style=verilog]
Insertion_block #(.W(W-1), .N(R_Q)) Ready_Queue(
    .data_in(data_in[W-2:0]),
    .data_out(task_out),
    .empty(empty),
    // ...
);
\end{lstlisting}

\textbf{Critical Queue (variable size):} Stores security-critical tasks requiring replica execution
\begin{lstlisting}[style=verilog]
Insertion_block #(.W(W-1), .N(16)) Critical_Queue(
    .data_in(data_in[W-2:0]),
    .data_out(task_critical_out),
    .empty(empty_critical),
    // ...
);
\end{lstlisting}

\textbf{Queue Selection Logic:}

When both queues contain tasks, the critical queue receives priority if its earliest task has a tighter effective deadline:

\begin{equation}
D_{eff,critical} = D_{relative} - E_{execution}
\end{equation}

\begin{lstlisting}[style=verilog, caption={Critical Task Priority Logic}]
assign check_critical = (~empty_critical) ?
    ((~empty) ?
        ((task_critical_out[31:16] - task_critical_out[15:0])
         <= task_out[31:16])
        : 1'b1)
    : 1'b0;
\end{lstlisting}

\subsubsection{Situation 1: Idle Processor Assignment}

When a processor completes its task, its deadline counter reaches zero, triggering a request signal. The arbiter selects one requesting processor per cycle.

\begin{lstlisting}[style=verilog, caption={Situation 1 Detection}]
// Check for any processor with deadline = 0
for (o = 0; o < CORE; o = o + 1) begin
    assign check_s1[o] = ~(|RT_in[o][31:16]) & RT_in[o][W-1];
end

Arbiter #(.CORE(CORE)) arbiter_tree(
    .req(check_s1),
    .grant_index(index_s1),
    .valid(find_s1)
);
\end{lstlisting}

When \texttt{find\_s1} is asserted, the earliest deadline task from either ready or critical queue is assigned to processor \texttt{index\_s1}.

\subsubsection{Situation 2: Preemptive Scheduling}

To prevent priority inversion, the scheduler iteratively checks all running tasks against newly arrived tasks. If a new task has an earlier deadline, it preempts the running task.

\begin{lstlisting}[style=verilog, caption={Situation 2 Preemption Logic}]
wire [15:0] tmp_s2 = check_critical ?
    (task_critical_out[31:16] - task_critical_out[15:0]) :
    task_out[31:16];

assign find_s2 = (
    (RT_before_S[s2_index][40]) ?  // Running task is critical
        tmp_s2 < (RT_before_S[s2_index][31:16] -
                  RT_before_S[s2_index][15:0]) :
        tmp_s2 < RT_before_S[s2_index][31:16]
    ) & (~empty | ~empty_critical) & CTRL_action;
\end{lstlisting}

The preempted task is written back to the appropriate queue:
\begin{itemize}
    \item Critical tasks (bit [40] = 1) → Critical Queue
    \item Normal tasks (bit [40] = 0) → Ready Queue
\end{itemize}

\textbf{Iteration Strategy:}

The \texttt{s2\_index} counter increments each cycle during the SCHEDULING state, checking one processor per cycle. Over 64 cycles, all 16 processors are checked multiple times, ensuring comprehensive preemption analysis.

\subsection{Module 4\_Controller.v - System Timing Controller}

The controller implements a deterministic state machine governing all timing aspects of SAFAS.

\begin{lstlisting}[style=verilog, caption={Controller State Machine}]
localparam SCHEDULING = 16'b0000_0000_0100_0000;  // 64
localparam IDLE = 16'b0000_0001_0000_0000;        // 256 (test)
                                                   // 65536 (prod)

always @(posedge clk) begin
    if (rst) begin
        state <= 3'b000;
        count <= 16'b0;
    end
    else begin
        case (state)
            3'b000: begin  // SCHEDULING
                if (count == SCHEDULING - 1) begin
                    state <= 3'b001;
                    count <= 16'b0;
                end
                else
                    count <= count + 1;
            end
            3'b001: begin  // MARGIN-TIME
                if (count == R_Q - 1) begin
                    state <= 3'b010;
                    count <= 16'b0;
                end
                else
                    count <= count + 1;
            end
            // ... additional states
        endcase
    end
end
\end{lstlisting}

\textbf{Control Signal Generation:}

\begin{table}[H]
\centering
\begin{tabular}{|l|l|p{6cm}|}
\hline
\textbf{Signal} & \textbf{Active State} & \textbf{Purpose} \\
\hline
action & SCHEDULING & Enables Situation 1 and Situation 2 checks \\
subtract & SUBTRACT & Single-cycle pulse for deadline/execution decrement \\
repair\_period & IDLE (partial) & Activates deadline miss detection and task removal \\
MQ\_active & SCHEDULING & Enables Main Queue to transfer tasks to scheduler \\
\hline
\end{tabular}
\caption{Controller Output Signals}
\end{table}

\subsection{Module 4\_Scheduler\_main.v - Top-Level Integration}

This module instantiates and connects all major components into a unified scheduler system.

\begin{lstlisting}[style=verilog, caption={Top-Level Module Instantiation}]
module Scheduler_main #(
    parameter W = 42,
    parameter R_Q = 64,
    parameter CORE = 16
)(
    input clk,
    input rst,
    input [W-1:0] task_in,
    input wr,
    output [W-1:0] task_exch,
    output v_exch
);

// Control Unit
Controller #(.R_Q(R_Q)) Control(
    .clk(clk),
    .rst(rst),
    .action(action),
    .subtract(subtract),
    .repair_period(repair_period),
    .MQ_active(MQ_active)
);

// Main Scheduler
Scheduler_partial #(.W(W), .R_Q(R_Q), .CORE(CORE)) Scheduler(
    .clk(clk),
    .rst(rst),
    .data_in(task_in),
    .wr(wr),
    // ... connections
    .RT_in(running_tasks),
    .RT_out(running_tasks_out)
);

// Processor Cores (simulation only)
Cores #(.W(W), .CORE(CORE)) Cores(
    .clk(clk),
    .rst(rst),
    .In(running_tasks_out),
    .Out(running_tasks),
    .subtract_en(subtract)
);

endmodule
\end{lstlisting}

\section{Verification Methodology}

\subsection{Testbench Architecture}

Three levels of testbenches verify SAFAS functionality:

\subsubsection{TB\_Insertion\_block.v - Unit Test}

Verifies basic EDF sorting with a 4-cell insertion block:

\begin{lstlisting}[style=verilog, caption={Insertion Block Test Sequence}]
initial begin
    // Initialize
    clk = 0;
    rst = 1;
    #10 rst = 0;

    // Insert tasks with varying deadlines
    data_in = 41'b10000000100000000000000110000000000000001;
    wr = 1; rd = 0;
    #10;

    data_in = 41'b10000000110000000000001000000000000000001;
    #10;

    // ... additional insertions

    // Enable subtraction
    subtract_en = 1;
    #10;
    subtract_en = 0;

    // Enable repair period
    repair_period = 1;
    #10;

    // Read out all tasks
    wr = 0; rd = 1;
    repeat (4) #10;

    $finish;
end
\end{lstlisting}

\textbf{Verification Points:}
\begin{itemize}
    \item Tasks emerge in ascending deadline order
    \item Overflow detection when queue is full
    \item Empty signal asserts correctly
    \item Deadline-missed tasks are invalidated during repair
\end{itemize}

\subsubsection{TB\_scheduler\_partial.v - Integration Test}

Tests scheduler with 8 processors and random task injection:

\begin{itemize}
    \item Random task generation using \texttt{\$random}
    \item Continuous write/read cycling
    \item Verification of task distribution across processors
    \item Monitoring of preemption events
\end{itemize}

\subsubsection{TB\_scheduler\_main.v - System-Level Validation}

The most comprehensive testbench simulates complete system operation with realistic workloads:

\textbf{Test Infrastructure:}

\begin{lstlisting}[style=verilog, caption={Main Testbench File I/O}]
localparam Size = 3275;  // Number of tasks in file
reg [W-1:0] buffer[0:Size-1];

initial begin
    $readmemb("C:/SAFAS/Experimental/TaskSet.txt", buffer);

    // Task injection during scheduling periods
    always @(posedge clk) begin
        if (CTRL_MQ_active && !MQ_empty && counter < Size) begin
            task_in <= buffer[counter];
            wr_MQ <= 1;
            counter <= counter + 1;
        end
        else
            wr_MQ <= 0;
    end
end
\end{lstlisting}

\textbf{Main Queue Instantiation:}

\begin{lstlisting}[style=verilog, caption={Main Queue (512 cells)}]
Insertion_block #(.W(W), .N(M_Q)) Main_Queue(
    .clk(clk),
    .rst(rst),
    .data_in(task_in_MQ),
    .wr(wr),
    .rd(rd),
    .data_out(task_to_scheduler),
    .empty(MQ_empty),
    .fail(v_exch),
    .data_fail(task_exch)
);
\end{lstlisting}

\textbf{Statistical Collection:}

The testbench collects comprehensive statistics from multiple sources:

\begin{lstlisting}[style=verilog, caption={Statistical Monitoring}]
// Aggregate deadline misses from all queues
always @(posedge clk) begin
    if (CTRL_repair_period) begin
        report_miss_tmp1 <= Ready_Q.report_miss;      // Ready Queue
        report_miss_tmp2 <= Main_Q.report_miss;       // Main Queue
        report_miss_tmp3 <= Critical_Q.report_miss;   // Critical Queue

        // Count misses
        RP_MISS_total <= RP_MISS_total +
            count_ones(report_miss_tmp1) +
            count_ones(report_miss_tmp2) +
            count_ones(report_miss_tmp3);
    end
end
\end{lstlisting}

\textbf{Simulation Termination:}

The simulation concludes when all tasks from the input file have been processed and all processors are idle:

\begin{lstlisting}[style=verilog, caption={Termination Condition}]
always @(posedge clk) begin
    if ((counter == Size) && (Cores.RP_current_tasks == 0)
        && MQ_empty) begin

        $display("\nFinal data analysis:");
        $display("SYSTEM:");
        $display("   All tasks: %d", RP_All_tasks);
        $display("   Finished tasks: %d", Cores.RP_PASS);
        $display("   Missed tasks: %d",
                 Cores.RP_MISS + RP_MISS_total);

        $display("\nSCHEDULER:");
        $display("   Received tasks: %d",
                 Main_Scheduler.RP_RECEIVED);
        $display("   Exchanged tasks: %d",
                 Main_Scheduler.RP_Exchanged);
        $display("   Preempted tasks: %d",
                 Main_Scheduler.RP_preempted);

        $finish;
    end
end
\end{lstlisting}

\subsection{Simulation Results}

Typical simulation results demonstrate SAFAS effectiveness:

\begin{table}[H]
\centering
\begin{tabular}{|l|r|}
\hline
\textbf{Metric} & \textbf{Value} \\
\hline
Total tasks released & 3,275 \\
Successfully completed tasks & 3,180 (97.1\%) \\
Deadline-missed tasks & 95 (2.9\%) \\
Tasks preempted & 142 \\
Queue overflow rejections & 12 \\
Scheduling iterations & 485 \\
Total clock cycles & 31,825,920 \\
Simulated time (@ 200 MHz) & 159.1 ms \\
\hline
\end{tabular}
\caption{Typical SAFAS Simulation Results}
\label{tab:sim_results}
\end{table}

\section{Task Generation and Workload Analysis}

\subsection{C++ Task Generator}

The \texttt{task\_generator} C++ program creates realistic real-time task workloads for testing SAFAS. It ensures tasks are schedulable under EDF constraints.

\subsubsection{Configuration Parameters}

\begin{lstlisting}[language=C++, caption={Task Generator Configuration}]
const int Num_cores = 16;           // Processor cores
float utilization = 14.5;           // Target utilization (0-16)
int Num_sec_tasks = 15;             // Security-critical tasks
float U_max = 0.3;                  // Max utilization per task
float margin = 0.025;               // Safety margin
int Num_replicas = ceil(3*(1-U_max)) + 1;  // Replicas
float wait_delay = 0.5;             // Time quantum (ms)
\end{lstlisting}

\subsubsection{EDF Feasibility Analysis}

For a set of tasks to be schedulable under EDF on $m$ processors:

\begin{equation}
\sum_{i=1}^{n} \frac{e_i}{d_i} \leq m
\end{equation}

where $e_i$ is execution time and $d_i$ is relative deadline of task $i$.

For security-critical tasks requiring $r$ replicas:

\begin{equation}
\sum_{i \in Critical} r \cdot \frac{e_i}{d_i} + \sum_{i \in Normal} \frac{e_i}{d_i} \leq m
\end{equation}

\textbf{Implementation:}

\begin{lstlisting}[language=C++, caption={Feasibility Check Algorithm}]
// Phase 1: Critical task feasibility
float U_critical = Num_cores * (1 - U_max) + U_max;

for (int i = 0; i < Num_sec_tasks; i++) {
    if (U_critical >= Num_replicas * tasks[i].u) {
        U_critical -= Num_replicas * tasks[i].u;
        tasks[i].feasible = true;
    }
}

// Phase 2: Normal task feasibility
float U_remaining = Num_cores - U_critical_used;

for (int i = Num_sec_tasks; i < NUM_TASK; i++) {
    if (U_remaining >= tasks[i].u) {
        U_remaining -= tasks[i].u;
        tasks[i].feasible = true;
    }
}
\end{lstlisting}

\subsubsection{Task Release Schedule}

Tasks are released periodically based on their deadline. For a simulation window [1000, 1400] time units:

\begin{lstlisting}[language=C++, caption={Periodic Task Release}]
for (int time = 1000; time <= 1400; time++) {
    for (int i = 0; i < NUM_TASK; i++) {
        if (tasks[i].feasible && (time % tasks[i].di == 0)) {
            int replicas = tasks[i].se ? Num_replicas : 1;

            for (int r = 0; r < replicas; r++) {
                // Output 42-bit task
                output_binary_task(tasks[i]);
            }
        }
    }
    output << "0" << endl;  // Separator
}
\end{lstlisting}

\subsection{Workload Characteristics}

Typical generated workloads exhibit:

\begin{table}[H]
\centering
\begin{tabular}{|l|r|}
\hline
\textbf{Characteristic} & \textbf{Value} \\
\hline
Total unique tasks & 79 \\
Security-critical tasks & 15 (19\%) \\
Normal tasks & 64 (81\%) \\
Replica factor for critical tasks & 3-4 \\
Average task utilization & 0.183 (18.3\%) \\
Total system utilization & 14.5 / 16 = 90.6\% \\
Deadline range & 4-256 time units \\
Tasks per release period & 8-12 \\
\hline
\end{tabular}
\caption{Typical Workload Characteristics}
\end{table}

\section{Security Features}

\subsection{Security-Critical Task Handling}

SAFAS provides specialized support for security-critical operations through:

\subsubsection{Replica Execution}

Critical tasks are replicated $r$ times to provide:
\begin{itemize}
    \item \textbf{Fault Tolerance:} Majority voting detects single-fault errors
    \item \textbf{Timing Attack Resistance:} Multiple executions complicate side-channel analysis
    \item \textbf{Redundancy:} System continues operating if some replicas fail
\end{itemize}

Replica count calculation:
\begin{equation}
r = \lceil 3 \cdot (1 - U_{max}) \rceil + 1
\end{equation}

For $U_{max} = 0.3$: $r = \lceil 2.1 \rceil + 1 = 4$ replicas

\subsubsection{Separate Critical Queue}

The dedicated critical task queue ensures:
\begin{itemize}
    \item Priority access when effective deadlines are tighter
    \item Isolation from normal task queue overflow
    \item Guaranteed scheduling attention
\end{itemize}

\subsubsection{Schedule-Based Attack Prevention}

Hardware scheduling inherently provides security advantages:
\begin{itemize}
    \item \textbf{Deterministic Timing:} Hardware execution eliminates cache timing variations
    \item \textbf{OS Isolation:} Scheduling logic separated from potentially compromised OS
    \item \textbf{Information Hiding:} Task scheduling state not visible to software
    \item \textbf{Preemption Control:} Hardware-enforced context switching
\end{itemize}

\subsection{Fault Detection}

Multiple mechanisms detect and handle failures:

\subsubsection{Deadline Miss Detection}

During the repair period, tasks with $D_{relative} < E_{execution}$ are flagged:

\begin{equation}
\text{missed} = (D_{relative} - E_{execution} < 0)
\end{equation}

These tasks are invalidated and statistics are recorded.

\subsubsection{Queue Overflow Protection}

When insertion blocks are full, overflow tasks are:
\begin{itemize}
    \item Rejected via \texttt{fail} signal
    \item Returned to OS via \texttt{task\_exch} output
    \item Counted in \texttt{RP\_Exchanged} statistics
\end{itemize}

\subsubsection{Statistical Monitoring}

Comprehensive counters track:
\begin{itemize}
    \item Tasks received, dispatched, completed
    \item Preemption events
    \item Deadline misses (per queue and global)
    \item Queue overflows
\end{itemize}

This provides visibility for system health monitoring and debugging.

\section{Performance Analysis}

\subsection{Hardware Resource Utilization}

Estimated resource usage for Xilinx Virtex FPGA implementation:

\begin{table}[H]
\centering
\begin{tabular}{|l|r|l|}
\hline
\textbf{Resource} & \textbf{Quantity} & \textbf{Notes} \\
\hline
\multicolumn{3}{|c|}{\textbf{Storage Elements}} \\
\hline
Ready Queue (64 cells × 41 bits) & 2,624 FFs & Main ready queue \\
Critical Queue (8-16 cells × 41 bits) & 328-656 FFs & Security tasks \\
Main Queue (512 cells × 42 bits) & 21,504 FFs & Input buffer \\
Running Tasks (16 × 42 bits) & 672 FFs & Processor state \\
Control registers & $\sim$100 FFs & State machine \\
\textbf{Total Flip-Flops} & \textbf{$\sim$25,200} & \\
\hline
\multicolumn{3}{|c|}{\textbf{Combinational Logic}} \\
\hline
Insertion cell comparators & 584 & 16-bit comparators \\
Arbiter tree & 31 & 2:1 arbiters \\
Multiplexers & $\sim$1,200 & Data routing \\
Arithmetic (subtract) & 584 & 16-bit decrementers \\
\hline
\multicolumn{3}{|c|}{\textbf{Estimated FPGA Usage}} \\
\hline
Slice Registers & $\sim$25,200 & 8-13\% of Virtex-6 LX240T \\
Slice LUTs & $\sim$15,000 & 10-15\% of Virtex-6 LX240T \\
Block RAMs & 0 & All distributed storage \\
DSP48 Slices & 0 & No multiplication \\
\hline
\end{tabular}
\caption{Estimated FPGA Resource Utilization}
\label{tab:resources}
\end{table}

\subsection{Timing Performance}

\subsubsection{Critical Path Analysis}

The longest combinational path determines maximum clock frequency:

\textbf{Insertion Cell Critical Path:}
\begin{enumerate}
    \item Register output (clock-to-Q delay): $\sim$0.5 ns
    \item 16-bit comparator: $\sim$2.0 ns
    \item 41-bit multiplexer: $\sim$1.5 ns
    \item Register setup time: $\sim$0.3 ns
\end{enumerate}

\textbf{Total Critical Path:} $\sim$4.3 ns $\rightarrow$ $f_{max} \approx 230$ MHz

\subsubsection{Scheduling Latency}

From task arrival to processor assignment:

\begin{table}[H]
\centering
\begin{tabular}{|l|r|l|}
\hline
\textbf{Stage} & \textbf{Cycles} & \textbf{Time @ 200 MHz} \\
\hline
Task written to Main Queue & 1 & 5 ns \\
Wait for SCHEDULING state & 0-65,536 & 0-328 $\mu$s \\
Transfer to Ready/Critical Queue & 1-64 & 5-320 ns \\
EDF sorting (pipeline fill) & 64 & 320 ns \\
Situation 1 arbitration & 4 & 20 ns \\
Task dispatch & 1 & 5 ns \\
\hline
\textbf{Average Latency} & $\sim$32,840 & $\sim$164 $\mu$s \\
\textbf{Maximum Latency} & $\sim$65,670 & $\sim$328 $\mu$s \\
\hline
\end{tabular}
\caption{Scheduling Latency Breakdown}
\end{table}

\subsubsection{Throughput}

\textbf{Task Processing Rate:}
\begin{itemize}
    \item Scheduling cycles per period: 64
    \item Tasks dispatched per cycle: 1 (Situation 1) + 0-1 (Situation 2)
    \item Maximum throughput: $\sim$2 tasks per 64 cycles = 1 task / 32 cycles
\end{itemize}

At 200 MHz: $\frac{200 \times 10^6}{32} \approx 6.25$ million tasks/second

\subsection{Scalability Analysis}

\subsubsection{Processor Count Scaling}

\begin{table}[H]
\centering
\begin{tabular}{|r|r|r|r|}
\hline
\textbf{Cores} & \textbf{Arbiter Levels} & \textbf{Registers} & \textbf{Latency} \\
\hline
4 & 2 & $\sim$6,500 & 2 cycles \\
8 & 3 & $\sim$13,000 & 3 cycles \\
16 & 4 & $\sim$25,200 & 4 cycles \\
32 & 5 & $\sim$50,000 & 5 cycles \\
64 & 6 & $\sim$100,000 & 6 cycles \\
\hline
\end{tabular}
\caption{Scaling with Processor Count}
\end{table}

\textbf{Observations:}
\begin{itemize}
    \item Storage grows linearly: $O(n)$
    \item Arbiter complexity grows logarithmically: $O(\log_2 n)$
    \item Critical path remains constant
\end{itemize}

\subsubsection{Queue Depth Scaling}

\begin{table}[H]
\centering
\begin{tabular}{|r|r|r|}
\hline
\textbf{Queue Depth} & \textbf{Flip-Flops} & \textbf{Insertion Time} \\
\hline
32 & $\sim$1,312 & 32 cycles \\
64 & $\sim$2,624 & 64 cycles \\
128 & $\sim$5,248 & 128 cycles \\
256 & $\sim$10,496 & 256 cycles \\
\hline
\end{tabular}
\caption{Queue Depth Impact}
\end{table}

Deeper queues increase buffering but proportionally increase sorting latency.

\section{Comparison with Software Schedulers}

\begin{table}[H]
\centering
\begin{tabular}{|l|c|c|}
\hline
\textbf{Characteristic} & \textbf{Software (OS)} & \textbf{SAFAS Hardware} \\
\hline
Scheduling Decision Time & 10-100 $\mu$s & 0.32 $\mu$s \\
Determinism & Variable & Cycle-accurate \\
CPU Overhead & 5-10\% & 0\% (dedicated HW) \\
Security Isolation & Vulnerable & Hardware-enforced \\
Maximum Cores Supported & 64-128 & Scalable (16-64) \\
EDF Sorting Complexity & $O(n \log n)$ & $O(n)$ pipeline \\
Context Switch Overhead & 5-20 $\mu$s & Hardware-direct \\
Power Efficiency & Moderate & High (specialized) \\
\hline
\end{tabular}
\caption{Software vs. Hardware Scheduler Comparison}
\end{table}

\section{Limitations and Future Work}

\subsection{Current Limitations}

\begin{enumerate}
    \item \textbf{Fixed Queue Sizes:} Queue depths are parameterized but must be determined at synthesis time
    \item \textbf{No Dynamic Priority:} EDF only; no support for priority-based or mixed scheduling
    \item \textbf{Limited Task Complexity:} Each task is assumed atomic (no internal dependencies)
    \item \textbf{No Task Migration:} Once assigned, tasks cannot move between processors
    \item \textbf{Simplified Processor Model:} Simulation assumes constant execution rates
\end{enumerate}

\subsection{Future Enhancements}

\subsubsection{Adaptive Queue Management}
Implement dynamic queue resizing or hierarchical queue structures to handle varying workload intensities.

\subsubsection{Multiple Scheduling Algorithms}
Add configurable scheduling policies:
\begin{itemize}
    \item Rate Monotonic Scheduling (RMS)
    \item Least Laxity First (LLF)
    \item Mixed-criticality support
\end{itemize}

\subsubsection{Task Dependency Handling}
Extend task format to include:
\begin{itemize}
    \item Predecessor/successor relationships
    \item Resource requirements
    \item Inter-task communication
\end{itemize}

\subsubsection{Power Management}
Implement Dynamic Voltage and Frequency Scaling (DVFS) integration:
\begin{itemize}
    \item Adjust processor speed based on workload
    \item Energy-aware scheduling decisions
    \item Thermal management support
\end{itemize}

\subsubsection{Enhanced Security}
\begin{itemize}
    \item Encrypted task communication channels
    \item Hardware-enforced memory isolation
    \item Formal verification of security properties
    \item Side-channel attack mitigation
\end{itemize}

\subsubsection{Machine Learning Integration}
Incorporate ML-based workload prediction to:
\begin{itemize}
    \item Anticipate task arrivals
    \item Optimize preemption decisions
    \item Adapt to changing workload patterns
\end{itemize}

\section{Conclusion}

SAFAS demonstrates the viability and advantages of hardware-based task scheduling for modern multi-core embedded systems. Through careful architectural design and efficient Verilog implementation, the system achieves:

\begin{itemize}
    \item \textbf{High Performance:} Sub-microsecond scheduling decisions with zero CPU overhead
    \item \textbf{Strong Determinism:} Cycle-accurate, predictable behavior critical for hard real-time systems
    \item \textbf{Security Awareness:} Hardware-enforced isolation and support for replicated critical task execution
    \item \textbf{Resource Efficiency:} Modest FPGA resource utilization (8-15\% of mid-range Virtex devices)
    \item \textbf{Scalability:} Logarithmic arbitration complexity enables scaling to 64+ cores
\end{itemize}

The verification methodology, including comprehensive testbenches and realistic workload generation, provides confidence in the design's correctness. Simulation results demonstrate successful scheduling of over 3,000 tasks with 97\% deadline satisfaction under high utilization (90\%).

As embedded systems continue to proliferate in safety-critical and security-sensitive applications—from autonomous vehicles to medical devices to industrial control systems—hardware schedulers like SAFAS offer a compelling path toward achieving the stringent timing, security, and reliability requirements that traditional software approaches struggle to provide.

The modular, parameterized Verilog implementation facilitates customization for specific application domains, while the comprehensive documentation and testbenches enable continued research and development. SAFAS serves as both a practical solution for real-time scheduling challenges and a foundation for future innovations in hardware-software co-design for embedded systems.

\section*{Acknowledgments}

This work builds upon research published in:

\begin{itemize}
    \item A. Norollah, H. Beitollahi, Z. Kazemi, and M. Fazeli, "A security-aware hardware scheduler for modern multi-core systems with hard real-time constraints," \textit{Elsevier Microprocessors and Microsystems (MICPRO)}, 2022. DOI: 10.1016/j.micpro.2022.104716.

    \item A. Norollah, Z. Kazemi, D. Derafshi, H. Beitollahi and M. Fazeli, "Protecting Security-Critical Real-Time Systems against Fault Attacks in Many-Core Platforms," \textit{2022 CPSSI 4th International Symposium on Real-Time and Embedded Systems and Technologies (RTEST)}, 2022, pp. 1-6. DOI: 10.1109/RTEST56034.2022.9850010.
\end{itemize}

\appendix

\section{Task Format Specification}

\subsection{Standard Task (41 bits)}
\begin{verbatim}
[40]: VALID (1=valid, 0=empty)
[39:32]: Task ID (8 bits, range 0-255)
[31:16]: Relative Deadline (16 bits, time units)
[15:0]: Execution Time (16 bits, time units)
\end{verbatim}

\subsection{Extended Task (42 bits)}
\begin{verbatim}
[41]: VALID (1=valid, 0=empty)
[40]: TYPE (1=security-critical, 0=normal)
[39:32]: Task ID (8 bits, range 0-255)
[31:16]: Relative Deadline (16 bits, time units)
[15:0]: Execution Time (16 bits, time units)
\end{verbatim}

\section{Module Hierarchy}

\begin{verbatim}
Scheduler_main (top-level)
├── Controller (4_Controller.v)
├── Scheduler_partial (3_Scheduler_partial.v)
│   ├── Insertion_block (Ready Queue, 1_Insertion_block.v)
│   │   └── Insertion_cell × 64 (1_Insertion_cell.v)
│   │       └── sub_task
│   ├── Insertion_block (Critical Queue, 1_Insertion_block.v)
│   │   └── Insertion_cell × 8-16
│   │       └── sub_task
│   └── Arbiter (2_Arbiter.v)
│       ├── Arbiter_cell_1 × 8 (Level 1)
│       ├── Arbiter_cell × 4 (Level 2)
│       ├── Arbiter_cell × 2 (Level 3)
│       └── Arbiter_cell × 1 (Level 4)
└── Cores (4_processors.v, simulation only)
    └── REG × 16 (0_REG.v)
\end{verbatim}

\section{Signal Reference}

\begin{table}[H]
\centering
\small
\begin{tabular}{|l|l|p{6cm}|}
\hline
\textbf{Signal Name} & \textbf{Width} & \textbf{Description} \\
\hline
clk & 1 & System clock (rising edge triggered) \\
rst & 1 & Synchronous active-high reset \\
task\_in & W & Input task from RTOS \\
wr & 1 & Write enable for task\_in \\
task\_exch & W & Exchanged/rejected task output \\
v\_exch & 1 & Valid exchange signal \\
action & 1 & Scheduling enabled signal \\
subtract & 1 & Deadline/execution decrement pulse \\
repair\_period & 1 & Deadline miss detection enable \\
MQ\_active & 1 & Main Queue active signal \\
RT\_in & W×CORE & Running tasks input array \\
RT\_out & W×CORE & Running tasks output array \\
empty & 1 & Queue empty indicator \\
fail & 1 & Queue overflow/rejection \\
check\_s1 & CORE & Situation 1 request signals \\
find\_s1 & 1 & Situation 1 valid signal \\
index\_s1 & $\log_2(CORE)$ & Selected processor index \\
find\_s2 & 1 & Situation 2 preemption signal \\
s2\_index & $\log_2(CORE)$ & Preemption target processor \\
\hline
\end{tabular}
\caption{Key Signal Descriptions}
\end{table}

\section{Parameter Reference}

\begin{table}[H]
\centering
\begin{tabular}{|l|r|l|}
\hline
\textbf{Parameter} & \textbf{Default} & \textbf{Description} \\
\hline
W & 41-42 & Task width in bits \\
N & 64 & Insertion block queue depth \\
R\_Q & 64 & Ready queue depth \\
M\_Q & 512 & Main queue depth \\
CORE & 16 & Number of processor cores \\
SCHEDULING & 64 & Scheduling state duration \\
IDLE & 65,536 & Idle state duration (production) \\
& 256 & Idle state duration (simulation) \\
\hline
\end{tabular}
\caption{Parameterization Options}
\end{table}

\end{document}
